\documentclass[a4paper,11pt]{article}
\usepackage[utf8]{inputenc}
\usepackage[T1]{fontenc}
\usepackage[ngerman]{babel}
\usepackage{geometry}
\geometry{left=2cm, right=2cm, top=2.5cm, bottom=2.5cm}
\usepackage{longtable}
\usepackage{booktabs}
\usepackage{array}
\usepackage{xcolor}
\usepackage{colortbl}
\usepackage{hyperref}
\usepackage{fancyhdr}
\usepackage{tabularx}
\usepackage{enumitem}

\definecolor{headerblue}{HTML}{1565C0}
\definecolor{tableheader}{HTML}{E3F2FD}
\definecolor{pass}{HTML}{C8E6C9}
\definecolor{inprog}{HTML}{FFF9C4}
\definecolor{todo}{HTML}{E0E0E0}

\hypersetup{colorlinks=true,linkcolor=headerblue,urlcolor=headerblue}

\pagestyle{fancy}
\fancyhf{}
\fancyhead[L]{\textbf{INSPECTAIR}}
\fancyhead[R]{Testspezifikation v2.0}
\fancyfoot[C]{\thepage}
\renewcommand{\headrulewidth}{0.4pt}

\setlength{\parindent}{0pt}
\setlength{\parskip}{6pt}

\begin{document}

\begin{center}
{\LARGE\textbf{INSPECTAIR}}\\[0.3cm]
{\Large Testspezifikation}\\[0.3cm]
{\large Testfälle und Abnahmekriterien}\\[0.5cm]
\begin{tabular}{ll}
\textbf{Projekt:} & InspectAir -- Luftqualitätsmessgerät \\
\textbf{Version:} & 2.0 \\
\textbf{Datum:} & Februar 2026 \\
\textbf{Autoren:} & Maximilian Liebl, Cedric Stroh, Jannis Messer \\
\textbf{Modul:} & Mikrocomputertechnik 1, DHBW Ravensburg \\
\textbf{Dozent:} & Thomas Stingl (Airbus) \\
\end{tabular}
\end{center}

\vspace{0.5cm}
\noindent\textit{Dieses Dokument wurde unter Verwendung von KI-Tools (Claude, Anthropic) zur Überprüfung erstellt.}
\vspace{0.3cm}

\tableofcontents
\newpage

% ============================================================
\section{Testübersicht}

\begin{tabularx}{\textwidth}{llXl}
\toprule
\rowcolor{tableheader}
\textbf{TC-ID} & \textbf{Kategorie} & \textbf{Testfall} & \textbf{Status} \\
\midrule
\rowcolor{pass} TC-01 & Einzeltest & AHT20 Temperatur/Feuchte & \checkmark~PASS \\
\rowcolor{pass} TC-02 & Einzeltest & SGP40 VOC-Sensor & \checkmark~PASS \\
\rowcolor{pass} TC-03 & Einzeltest & MH-Z19C CO\textsubscript{2}-Sensor & \checkmark~PASS \\
\rowcolor{pass} TC-04 & Einzeltest & PMS5003 Feinstaub & \checkmark~PASS \\
\rowcolor{pass} TC-05 & Einzeltest & LD2410C Radar & \checkmark~PASS \\
\rowcolor{pass} TC-06 & Einzeltest & TFT-Display (LVGL 9) & \checkmark~PASS \\
\rowcolor{pass} TC-07 & Hardware & BS170 MOSFET Backlight & \checkmark~PASS \\
\rowcolor{inprog} TC-08 & Hardware & Stützkondensatoren & $\circlearrowright$ \\
\rowcolor{inprog} TC-09 & Integration & Alle Sensoren parallel & $\circlearrowright$ \\
\rowcolor{todo} TC-10 & UI & Overview-Screen & $\circ$ \\
\rowcolor{todo} TC-11 & UI & Detail-Screen & $\circ$ \\
\rowcolor{todo} TC-12 & UI & Button-Navigation & $\circ$ \\
\rowcolor{todo} TC-13 & UI & Farbkonsistenz & $\circ$ \\
\rowcolor{todo} TC-14 & Power & Display-Dimming & $\circ$ \\
\rowcolor{todo} TC-15 & Power & Radar-Wake-Up & $\circ$ \\
\rowcolor{todo} TC-16 & Power & Light Sleep & $\circ$ \\
\rowcolor{todo} TC-17 & System & Gesamtszenario & $\circ$ \\
\bottomrule
\end{tabularx}

% ============================================================
\section{Einzeltests (TC-01 bis TC-06)}

\subsection{TC-01: AHT20 Temperatur- und Feuchtemessung}
\begin{tabularx}{\textwidth}{lX}
\toprule
\rowcolor{tableheader}
\textbf{Parameter} & \textbf{Wert} \\
\midrule
Anforderung & REQ-F-004, REQ-HW-005 \\
Vorbedingung & AHT20 an I2C (SDA:~GPIO8, SCL:~GPIO18), 4{,}7\,k$\Omega$ Pull-ups \\
Durchführung & \texttt{aht\_sgp\_init()} aufrufen, 10 Messzyklen mit \texttt{aht\_sgp\_read()} \\
Erwartetes Ergebnis & Temperatur: 15--30\,°C ($\pm$0{,}3\,°C), Feuchte: 20--80\,\% ($\pm$2\,\%) \\
Prüfung & Vergleich mit Referenz-Thermometer/-Hygrometer \\
Status & \checkmark~PASS -- Werte plausibel und stabil \\
\bottomrule
\end{tabularx}

\subsection{TC-02: SGP40 VOC-Messung}
\begin{tabularx}{\textwidth}{lX}
\toprule
\rowcolor{tableheader}
\textbf{Parameter} & \textbf{Wert} \\
\midrule
Anforderung & REQ-F-003, REQ-HW-005 \\
Vorbedingung & SGP40 am gleichen I2C-Bus wie AHT20, Addr: 0x59 \\
Durchführung & \texttt{aht\_sgp\_read()} mit VOC-Index-Algorithmus, Aufwärmzeit 60\,s \\
Erwartetes Ergebnis & VOC-Index: 0--500, Basiswert nach Aufwärmung 80--120 \\
Prüfung & Provokationstest: Alkohol/Desinfektionsmittel in Sensornähe $\to$ Index steigt \\
Status & \checkmark~PASS -- VOC-Index reagiert auf Ausdünstungen \\
\bottomrule
\end{tabularx}

\subsection{TC-03: MH-Z19C CO\textsubscript{2}-Messung}
\begin{tabularx}{\textwidth}{lX}
\toprule
\rowcolor{tableheader}
\textbf{Parameter} & \textbf{Wert} \\
\midrule
Anforderung & REQ-F-001, REQ-HW-003 \\
Vorbedingung & MH-Z19C über SoftwareSerial (RX:~GPIO4, TX:~GPIO5), 220\,\textmu F Kondensator, 3\,Min Aufwärmzeit \\
Durchführung & \texttt{mhz19c\_read()}, 10 Messzyklen im Abstand von 5\,s \\
Erwartetes Ergebnis & CO\textsubscript{2}: 400--2000\,ppm, \texttt{valid}-Flag = \texttt{true} \\
Prüfung & Ausatmen in Sensornähe $\to$ CO\textsubscript{2}-Wert steigt auf $>$1000\,ppm \\
Status & \checkmark~PASS -- Werte reagieren zuverlässig auf Atemluft \\
\bottomrule
\end{tabularx}

\subsection{TC-04: PMS5003 Feinstaubmessung}
\begin{tabularx}{\textwidth}{lX}
\toprule
\rowcolor{tableheader}
\textbf{Parameter} & \textbf{Wert} \\
\midrule
Anforderung & REQ-F-002, REQ-HW-004 \\
Vorbedingung & PMS5003 über UART1 (RX:~GPIO16, TX:~GPIO17), interner Lüfter aktiv \\
Durchführung & \texttt{pms5003\_read()}, 10 Messzyklen \\
Erwartetes Ergebnis & PM2.5: 0--300\,\textmu g/m³, PM10: 0--500\,\textmu g/m³ \\
Prüfung & Staubquelle (Kreidestaub) in Nähe $\to$ PM-Werte steigen \\
Status & \checkmark~PASS -- Reaktion auf Partikelquellen nachgewiesen \\
\bottomrule
\end{tabularx}

\subsection{TC-05: LD2410C Radar-Präsenzerkennung}
\begin{tabularx}{\textwidth}{lX}
\toprule
\rowcolor{tableheader}
\textbf{Parameter} & \textbf{Wert} \\
\midrule
Anforderung & REQ-F-005, REQ-HW-006 \\
Vorbedingung & LD2410C über UART2 (RX:~GPIO7, TX:~GPIO6), TXS0108E Level Shifter, 220\,\textmu F Kondensator \\
Durchführung & \texttt{ld2410c\_read()}, Person bewegt sich im Raum \\
Erwartetes Ergebnis & \texttt{presence}: true bei Person im Raum (bis 6\,m Reichweite), false bei leerem Raum \\
Prüfung & Raum betreten/verlassen, Entfernungswert prüfen \\
Status & \checkmark~PASS -- Zuverlässige Erkennung bis 5\,m \\
\bottomrule
\end{tabularx}

\subsection{TC-06: TFT-Display und LVGL 9}
\begin{tabularx}{\textwidth}{lX}
\toprule
\rowcolor{tableheader}
\textbf{Parameter} & \textbf{Wert} \\
\midrule
Anforderung & REQ-F-007, REQ-HW-007 \\
Vorbedingung & ST7796S über SPI, LovyanGFX 1.1.16, LVGL 9.2.2 \\
Durchführung & \texttt{lvgl\_init()} + \texttt{ui\_init()}, alle 5 Screens durchschalten \\
Erwartetes Ergebnis & Alle Screens werden korrekt gerendert, keine Artefakte, 480$\times$320 Auflösung \\
Prüfung & Visuelle Inspektion aller 5 Screens (Tree, Overview, Detail, Analog, Bubble) \\
Status & \checkmark~PASS -- Alle Screens rendering korrekt \\
\bottomrule
\end{tabularx}

% ============================================================
\section{Hardware-Tests (TC-07, TC-08)}

\subsection{TC-07: BS170 MOSFET Backlight-Steuerung}
\begin{tabularx}{\textwidth}{lX}
\toprule
\rowcolor{tableheader}
\textbf{Parameter} & \textbf{Wert} \\
\midrule
Anforderung & REQ-F-006 \\
Vorbedingung & BS170 MOSFET mit 10\,k$\Omega$ Gate-Widerstand an GPIO3 \\
Durchführung & \texttt{lvgl\_setBrightness()} mit Werten 0, 30, 128, 255 aufrufen \\
Erwartetes Ergebnis & Display-Helligkeit ändert sich stufenlos entsprechend PWM-Wert \\
Status & \checkmark~PASS \\
\bottomrule
\end{tabularx}

\subsection{TC-08: Stützkondensatoren Stabilität}
\begin{tabularx}{\textwidth}{lX}
\toprule
\rowcolor{tableheader}
\textbf{Parameter} & \textbf{Wert} \\
\midrule
Anforderung & REQ-HW-010 \\
Vorbedingung & 220\,\textmu F Elkos an VCC der MH-Z19C und LD2410C \\
Durchführung & 30\,Min Dauerbetrieb mit allen Sensoren gleichzeitig aktiv \\
Erwartetes Ergebnis & Kein Systemabsturz, keine Messwert-Aussetzer, stabile 5V/3{,}3V-Versorgung \\
Status & $\circlearrowright$~In Bearbeitung \\
\bottomrule
\end{tabularx}

% ============================================================
\section{Integrationstests (TC-09)}

\subsection{TC-09: Alle Sensoren parallel}
\begin{tabularx}{\textwidth}{lX}
\toprule
\rowcolor{tableheader}
\textbf{Parameter} & \textbf{Wert} \\
\midrule
Anforderung & REQ-F-001 bis REQ-F-005 \\
Vorbedingung & Alle 5 Sensoren verkabelt und einzeln getestet (TC-01 bis TC-05) \\
Durchführung & Hauptprogramm starten, 15\,Min Dauerbetrieb, alle Sensorwerte auf Serial Monitor prüfen \\
Erwartetes Ergebnis & Alle Sensoren liefern gleichzeitig plausible Werte, kein UART-Konflikt, I2C stabil \\
Status & $\circlearrowright$~In Bearbeitung \\
\bottomrule
\end{tabularx}

% ============================================================
\section{UI-Tests (TC-10 bis TC-13)}

\subsection{TC-10: Overview-Screen}
\begin{tabularx}{\textwidth}{lX}
\toprule
\rowcolor{tableheader}
\textbf{Parameter} & \textbf{Wert} \\
\midrule
Anforderung & REQ-F-007, REQ-F-009 \\
Durchführung & Overview-Screen aufrufen, Sensorwerte variieren \\
Erwartetes Ergebnis & Große AQI-Anzeige rechts, 2 Kacheln (Temperatur, Feuchte) links, Uhrzeit/Datum oben, Farbcodierung korrekt \\
Status & $\circ$~Offen \\
\bottomrule
\end{tabularx}

\subsection{TC-11: Detail-Screen}
\begin{tabularx}{\textwidth}{lX}
\toprule
\rowcolor{tableheader}
\textbf{Parameter} & \textbf{Wert} \\
\midrule
Anforderung & REQ-F-007, REQ-F-009 \\
Durchführung & Detail-Screen aufrufen, alle Messwerte prüfen \\
Erwartetes Ergebnis & Kleine AQI-Anzeige, 4 Kacheln mit CO\textsubscript{2}, PM2.5, VOC, Temperatur/Feuchte, korrekte Einheiten und Farbcodierung \\
Status & $\circ$~Offen \\
\bottomrule
\end{tabularx}

\subsection{TC-12: Button-Navigation}
\begin{tabularx}{\textwidth}{lX}
\toprule
\rowcolor{tableheader}
\textbf{Parameter} & \textbf{Wert} \\
\midrule
Anforderung & REQ-F-009 \\
Vorbedingung & UI-Button an GPIO1 (Active LOW, interner Pull-up) \\
Durchführung & Button wiederholt drücken, Screen-Wechsel prüfen \\
Erwartetes Ergebnis & Zyklische Navigation: Tree $\to$ Overview $\to$ Detail $\to$ Analog $\to$ Bubble $\to$ Tree. Entprellung: 250\,ms Debounce \\
Status & $\circ$~Offen \\
\bottomrule
\end{tabularx}

\subsection{TC-13: Farbkonsistenz über Grenzwerte}
\begin{tabularx}{\textwidth}{lX}
\toprule
\rowcolor{tableheader}
\textbf{Parameter} & \textbf{Wert} \\
\midrule
Anforderung & REQ-F-008 \\
Durchführung & Verschiedene Messwerte provozieren und Farbänderungen auf allen Screens prüfen \\
Erwartetes Ergebnis & Identische Farbcodierung auf allen 5 Screens gemäß Grenzwert-Tabelle (Anhang~A) \\
Status & $\circ$~Offen \\
\bottomrule
\end{tabularx}

% ============================================================
\section{Power-Management-Tests (TC-14 bis TC-16)}

\subsection{TC-14: Display-Dimming}
\begin{tabularx}{\textwidth}{lX}
\toprule
\rowcolor{tableheader}
\textbf{Parameter} & \textbf{Wert} \\
\midrule
Anforderung & REQ-F-006 \\
Durchführung & Raum verlassen, 45\,s warten \\
Erwartetes Ergebnis & Display dimmt von Brightness 255 auf 30 mit Fade-Übergang \\
Status & $\circ$~Offen \\
\bottomrule
\end{tabularx}

\subsection{TC-15: Radar-Wake-Up}
\begin{tabularx}{\textwidth}{lX}
\toprule
\rowcolor{tableheader}
\textbf{Parameter} & \textbf{Wert} \\
\midrule
Anforderung & REQ-F-005, REQ-F-006 \\
Durchführung & Display im DIMMED- oder OFF-Zustand, Raum betreten \\
Erwartetes Ergebnis & Radar erkennt Präsenz $\to$ \texttt{wakeUp()} $\to$ Display auf volle Helligkeit (255) \\
Status & $\circ$~Offen \\
\bottomrule
\end{tabularx}

\subsection{TC-16: Light Sleep}
\begin{tabularx}{\textwidth}{lX}
\toprule
\rowcolor{tableheader}
\textbf{Parameter} & \textbf{Wert} \\
\midrule
Anforderung & REQ-NF-003 \\
Durchführung & Raum für $>$5\,Min verlassen, Sleep-Eintritt prüfen \\
Erwartetes Ergebnis & ESP32 wechselt in Light Sleep ($\sim$0{,}8\,mA), Aufwachzeit $<$5\,ms, RAM-Daten erhalten \\
Status & $\circ$~Offen \\
\bottomrule
\end{tabularx}

% ============================================================
\section{System-Gesamttest (TC-17)}

\subsection{TC-17: Vollständiges Betriebsszenario}
\begin{tabularx}{\textwidth}{lX}
\toprule
\rowcolor{tableheader}
\textbf{Parameter} & \textbf{Wert} \\
\midrule
Anforderung & Alle REQ-F und REQ-NF \\
Durchführung & System einschalten, 60\,Min Dauerbetrieb: Alle Screens durchschalten, Grenzwerte provozieren, Raum betreten/verlassen, WiFi-Verbindung und NTP prüfen, Web Remote testen \\
Erwartetes Ergebnis & Stabile Messwerte, korrekte Farbcodierung, Power Management reagiert auf Präsenz, keine Abstürze, Uhrzeit korrekt \\
Status & $\circ$~Offen \\
\bottomrule
\end{tabularx}

% ============================================================
\section*{Anhang A: Grenzwert-Konfiguration}
\addcontentsline{toc}{section}{Anhang A: Grenzwert-Konfiguration}

\begin{tabularx}{\textwidth}{lllll}
\toprule
\rowcolor{tableheader}
\textbf{Parameter} & \textbf{Grün (Gut)} & \textbf{Gelb (Mäßig)} & \textbf{Orange (Schlecht)} & \textbf{Rot (Gefährlich)} \\
\midrule
CO\textsubscript{2} (ppm) & $<$800 & 800--1000 & 1000--1500 & $>$1500 \\
PM2.5 (\textmu g/m³) & $<$5 & 5--15 & 15--25 & $>$25 \\
PM10 (\textmu g/m³) & $<$15 & 15--45 & 45--75 & $>$75 \\
VOC-Index & $<$100 & 100--200 & 200--300 & $>$300 \\
Temperatur (°C) & 19--23 & 17--19 / 23--26 & 15--17 / 26--28 & $<$15 / $>$28 \\
Feuchte (\%) & 40--60 & 30--40 / 60--70 & 20--30 / 70--80 & $<$20 / $>$80 \\
\bottomrule
\end{tabularx}

\vspace{0.3cm}
\textit{Grenzwerte basieren auf WHO Air Quality Guidelines 2021 (PM2.5, PM10) und UBA-Empfehlungen (CO\textsubscript{2}, Temperatur, Feuchte). Implementierung in \texttt{colors.h}.}

% ============================================================
\section*{Anhang B: Timeout-Konfiguration}
\addcontentsline{toc}{section}{Anhang B: Timeout-Konfiguration}

\begin{tabularx}{\textwidth}{lll}
\toprule
\rowcolor{tableheader}
\textbf{Parameter} & \textbf{Wert} & \textbf{Quelle} \\
\midrule
PRESENCE\_TIMEOUT\_DIM\_MS & 45\,000\,ms (45\,s) & \texttt{power\_manager.h} \\
PRESENCE\_TIMEOUT\_OFF\_MS & 300\,000\,ms (5\,Min) & \texttt{power\_manager.h} \\
BRIGHTNESS\_ACTIVE & 255 (100\,\%) & \texttt{power\_manager.h} \\
BRIGHTNESS\_DIMMED & 30 ($\sim$12\,\%) & \texttt{power\_manager.h} \\
BRIGHTNESS\_OFF & 0 & \texttt{power\_manager.h} \\
Button Debounce & 250\,ms & \texttt{main.cpp} \\
Sensor Read Interval & 2\,000--5\,000\,ms & \texttt{main.cpp} \\
\bottomrule
\end{tabularx}

% ============================================================
\section*{Anhang C: Kritische Hardware-Hinweise}
\addcontentsline{toc}{section}{Anhang C: Kritische Hardware-Hinweise}

\textbf{WARNUNG: Vor dem Testen unbedingt beachten:}

\begin{enumerate}[nosep]
\item \textbf{Level Shifter PFLICHT:} LD2410C-TX sendet 5V-Signale. Ohne TXS0108E wird GPIO7 des ESP32 zerstört!
\item \textbf{Stützkondensatoren PFLICHT:} 220\,\textmu F Elkos an VCC von MH-Z19C und LD2410C. Ohne Kondensatoren kommt es zu Systemabstürzen.
\item \textbf{I2C Pull-ups prüfen:} AHT20 und SGP40 benötigen 4{,}7\,k$\Omega$ Pull-ups an SDA (GPIO8) und SCL (GPIO18).
\item \textbf{PMS5003 Belüftung:} Der interne Lüfter benötigt freien Luftstrom. Im Gehäuse Belüftungsöffnungen vorsehen.
\item \textbf{GPIO-Belegung:} GPIO0 nicht verwenden (Boot-Pin!), GPIO10--13 für Flash/PSRAM reserviert.
\end{enumerate}

% ============================================================
\section{Änderungshistorie}

\begin{tabularx}{\textwidth}{l l l X}
\toprule
\rowcolor{tableheader}
\textbf{Version} & \textbf{Datum} & \textbf{Autor} & \textbf{Änderung} \\
\midrule
1.0 & Januar 2026 & Team & Erstversion \\
2.0 & Februar 2026 & Team & Grenzwerte aktualisiert (WHO 2021, 4-stufig), Pin-Belegung aus \texttt{pins.h}, SoftwareSerial für MH-Z19C, 5 Screens statt 2, Power-Management-Tests ergänzt (TC-14 bis TC-16), Timeout-Konfiguration aus Code, Hardware-Warnungen ergänzt. \\
\bottomrule
\end{tabularx}

\end{document}
